\chapter{Problem \& Reflections}

Many of the hardest problems in this project were not purely visual and not purely logical. In practice, presentation, interaction, testing, and planning were tightly coupled. Once the interface became more elaborate, every UI change exposed another layer of engineering work: raycast boundaries, runtime layout assumptions, battle-state messaging, or proof that a fix had actually held. This chapter reflects on those issues at technical, project-management, and team-management levels.

\section{Technical Problem}

One recurring technical problem was that visual polish changed interaction behaviour. As the title overlay, themed HUD, and hand dock became more elaborate, previously hidden assumptions in the prototype started to fail. Concrete examples included title overlay content colliding vertically, the top-left prompt stack overlapping turn and energy labels, the hand dock obscuring lower units, and action availability being visible only in the battle log.

The important lesson was that these defects could not be handled with visual tweaking alone. Several bugs existed at the boundary between layout and logic. A panel could look harmless but still intercept clicks; a debug shortcut could speed up testing but bypass the same flow path used by the real HUD; a screenshot could show the board correctly while missing overlay elements entirely.

Another technical limitation came from tooling. Overlay screenshots were not always trustworthy for `Screen Space Overlay` verification in this setup. The practical response was to split validation responsibility: hierarchy assertions and runtime tests for structural proof, and manual smoke for final overlay readability.

\section{Project Management}

From a project-management perspective, the main challenge was resisting attractive but expensive scope growth. The game concept could easily have drifted toward a larger tactics RPG with more cards, more maps, more animation, and progression systems. Instead, the team kept returning to the vertical-slice rule: build one encounter well, then extend only when the current layer is stable.

This project-management choice is closely related to the development model described in Chapter 4. The team worked iteratively in a lightweight agile style, regularly returning to the backlog after each usable slice \cite{agileManifesto2001}. That proved more appropriate than trying to freeze all design decisions up front in a strict sequential process, especially once visible usability problems began to emerge. In retrospect, this also matches Royce's warning that large software work becomes risky when treated as a simple one-pass flow \cite{royce1970}.

Testing and evidence collection also became planning tools rather than only reporting tasks. When a feature area lacked a clear validation path, that was treated as a planning warning sign. Improvements such as disabled-card overlays, world-space health bars, and menu-flow fixes were prioritised because they closed real usability gaps and could be backed by tests or smoke evidence. Cosmetic work that could not yet be validated or stabilised was deprioritised.

\section{Team Management}

The team-management story is more mixed, and it is important to describe it honestly. Early coordination relied more on informal communication and shared files than on a fully preserved repository-first workflow. Later, the project converged onto a clearer branch-and-merge model with a stronger evidence tree. That change improved accountability and made the final codebase and report trail easier to audit, but it also meant that some early process context had to be preserved as recovered notes rather than as perfect historical records.

This was not only an administrative issue. It affected how the team divided work and reviewed risk. Once responsibilities, issue logs, screenshots, and report material lived in the same workspace, ambiguity reduced sharply. Code ownership became easier to follow, design changes became easier to justify, and late-stage fixes became easier to validate collaboratively.

The most positive team-management lesson is therefore standardisation. The project became stronger once the team treated one repository, one evidence tree, and one final-report source as the canonical workspace. The most important remaining weakness is that this standardisation happened later than ideal, so the report must still be transparent about which early records are formal and which are recovered.
